% !TEX program = xelatex
% Шаблон научной статьи НИТУ «МИСиС»
% Заполните блоки ниже и скомпилируйте: xelatex article.tex

\documentclass{misisarticle}

% --- Классификационные коды (оставьте пустым {}, если не требуются) ---
\grnti{28.23.15, 20.23.17, 28.23.37}
\asjc{1702, 1710, 1706}
\oecd{1.02}

% --- Название ---
\title{Сравнение pointwise, pairwise и listwise ранжирования новостей разными LLM на датасете MIND}

% --- Авторы: \misisauthor{ФИО}{аффилиация}{город, страна}{email} ---
\misisauthor{С.~О.~Галачинский}{кафедра инженерной кибернетики НИТУ «МИСиС»}{Москва, Россия}{student@edu.misis.ru}
\misisauthor{А.~С.~Лыкошин}{кафедра инженерной кибернетики НИТУ «МИСиС»}{Москва, Россия}{supervisor@misis.ru}

% --- Аннотация и ключевые слова ---
\abstract{%
  Рассматривается задача ранжирования новостных сообщений большими языковыми моделями (LLM) без дообучения (zero-shot). Известны три парадигмы промптинга: pointwise, pairwise и listwise. Однако их сравнение именно на новостных данных, где важны свежесть и разнообразие выдачи, проведено недостаточно, а выбор парадигмы определяет и качество, и стоимость вывода. Цель работы — определить, какая парадигма даёт лучший компромисс между качеством ранжирования и вычислительной стоимостью в новостном домене. Для этого реализуются три режима ранжирования для нескольких LLM (закрытая модель через API и открытая модель), результаты сравниваются с базовыми методами (BM25, эмбеддинги, нейросетевая модель NRMS). Проверка выполняется на датасете MIND на подвыборке показов по метрикам AUC, MRR, nDCG@5 и nDCG@10, а также по числу израсходованных токенов и времени работы.
}

\keywords{%
  ранжирование новостей, большие языковые модели, переранжирование, zero-shot промптинг, рекомендательные системы, датасет MIND
}

\begin{document}

\maketitle

% ============================================================
\section{Введение}

Системы персонализированных рекомендаций новостей играют ключевую роль в современных информационных платформах, помогая пользователям ориентироваться в ежедневном потоке публикаций. В отличие от классических областей рекомендаций, таких как электронная коммерция или стриминговые сервисы, новостной домен характеризуется коротким жизненным циклом объектов, постоянным появлением новых статей (проблема холодного старта) и быстрой сменой пользовательских интересов~\cite{wu2020mind}. Традиционные алгоритмы лексического поиска (например, BM25) не способны улавливать глубокие семантические связи между заголовками, а специализированные нейросетевые архитектуры, такие как NRMS~\cite{wu2019nrms}, требуют регулярного переобучения на больших объёмах размеченных логов кликов.

Появление больших языковых моделей (Large Language Models, LLM) открыло возможность решения задачи ранжирования в режиме без дополнительного обучения (\textit{zero-shot}), опираясь на заложенные в модель при предобучении знания о мире и способность следовать инструкциям~\cite{zhang2025recommendation, sosedka2024telegram}. Вместе с тем прямое применение LLM в качестве ранжирующих агентов сталкивается с необходимостью выбора стратегии формирования промпта. В литературе выделяют три основные парадигмы: поэлементную (\textit{pointwise}), попарную (\textit{pairwise}) и списочную (\textit{listwise})~\cite{sun2023chatgpt, qin2024pairwise}. Каждая из них накладывает собственные ограничения на длину контекста, количество обращений к модели, время отклика и итоговую точность упорядочивания.

Целью данной работы является исследование компромисса между качеством ранжирования и вычислительной стоимостью трёх парадигм \textit{zero-shot} промптинга для открытых и проприетарных LLM на материале крупномасштабного новостного датасета MIND~\cite{wu2020mind}.

% ============================================================
\section{Анализ работ}

Исследования в области персонализированного ранжирования новостей активно развиваются благодаря публикации эталонного датасета Microsoft News Dataset (MIND)~\cite{wu2020mind}. В качестве базовых специализированных решений широкое распространение получили модели с механизмами внутреннего внимания, в частности архитектура NRMS (Neural News Recommendation with Multi-Head Self-Attention)~\cite{wu2019nrms}, которая строит векторные представления новостей и профиля пользователя на основе истории просмотренных заголовков. Несмотря на высокую точность на фиксированной выборке, подобные модели требуют затратного обучения с учителем и теряют эффективность при резком изменении информационной повестки.

С развитием генеративных предобученных трансформеров фокус исследований сместился в сторону использования LLM в прикладных корпоративных и рекомендательных системах~\cite{beresnev2023chatbots, sosedka2024telegram}. В работе Zhang et al.~\cite{zhang2025recommendation} предложен подход InstructRec, в котором задача рекомендации формулируется как следование текстовой инструкции (\textit{instruction following}). Авторы показали, что универсальные языковые модели способны успешно оценивать релевантность кандидатов на основе текстового описания предпочтений пользователя в \textit{pointwise}-режиме.

В исследовании Sun et al.~\cite{sun2023chatgpt} продемонстрирована эффективность использования закрытых моделей семейства GPT в качестве агентов переранжирования (\textit{re-ranking agents}) в \textit{listwise}-формате. Модель получает на вход сразу список кандидатов и генерирует перестановку их идентификаторов в порядке убывания релевантности. Для обхода ограничений на размер контекстного окна авторы предложили стратегию скользящего окна (\textit{sliding window}).

В свою очередь, Qin et al.~\cite{qin2024pairwise} выявили, что \textit{listwise}-промптинг часто оказывается слишком сложным для моделей среднего размера и приводит к нарушению формата вывода. Авторы предложили парадигму попарного ранжирования (Pairwise Ranking Prompting, PRP), в которой LLM сравнивает только два документа за один запрос. Было показано, что \textit{pairwise}-подход обеспечивает высокое качество даже на открытых моделях умеренного размера, однако вопрос соотношения точности и вычислительных затрат всех трёх подходов применительно к специфике новостных лент остаётся недостаточно изученным.

% ============================================================
\section{Постановка задачи}

Пусть для некоторого пользователя $u$ задана история его предыдущих взаимодействий (кликов) $H_u$, состоящая из $m$ ранее прочитанных статей $d_1, \dots, d_m$, где каждый элемент характеризуется заголовком, категорией и кратким содержанием. В момент показа (сессии импрессии) пользователю предлагается набор из $N$ кандидатных новостей $C_u$, включающий элементы $c_1, \dots, c_N$. Для каждого кандидата известна истинная бинарная метка релевантности $y_i \in \{0, 1\}$, где $y_i = 1$ означает наличие клика пользователя по новости $c_i$, а $y_i = 0$~--- отсутствие клика.

Требуется построить ранжирующий алгоритм $f_{\mathrm{LLM}}(H_u, C_u)$ на базе большой языковой модели без дообучения весов (\textit{zero-shot}), который ставит в соответствие списку кандидатов $C_u$ упорядоченную перестановку индексов $\pi$ или вектор оценок релевантности $\hat{s}$ так, чтобы максимизировать метрики качества ранжирования при минимизации вычислительных ресурсов.

Для оценки качества ранжирования на подвыборке сессий показов датасета MIND~\cite{wu2020mind} используются следующие метрики (формулы~\eqref{eq:auc}--\eqref{eq:ndcg}):

\begin{itemize}
  \item площадь под кривой ошибок (\textbf{AUC}), оценивающая вероятность того, что случайно выбранная релевантная новость получит более высокую оценку, чем нерелевантная;
  \item средний обратный ранг первой релевантной новости в итоговом списке (\textbf{MRR});
  \item нормализованный дисконтированный кумулятивный выигрыш на срезе первых $K$ позиций (\textbf{nDCG@K}) для $K \in \{5, 10\}$.
\end{itemize}

Расчётные формулы метрик имеют вид:
\begin{equation}
  \mathrm{AUC} = \frac{1}{|P| \cdot |Q|} \sum_{i \in P} \sum_{j \in Q} \mathbb{I}(\hat{s}_i > \hat{s}_j),
  \label{eq:auc}
\end{equation}
\begin{equation}
  \mathrm{MRR} = \frac{1}{|U|} \sum_{u \in U} \frac{1}{\mathrm{rank}_u^*},
  \label{eq:mrr}
\end{equation}
\begin{equation}
  \mathrm{DCG}@K = \sum_{i=1}^{K} \frac{y_{\pi_i}}{\log_2(i + 1)},
  \label{eq:dcg}
\end{equation}
\begin{equation}
  \mathrm{nDCG}@K = \frac{\mathrm{DCG}@K}{\mathrm{IDCG}@K},
  \label{eq:ndcg}
\end{equation}
где $P$ и $Q$~--- множества позитивных ($y_i = 1$) и негативных ($y_j = 0$) кандидатов в показе, $\mathbb{I}(\cdot)$~--- индикаторная функция, $\mathrm{rank}_u^*$~--- позиция первой кликнутой новости в выдаче для показа $u \in U$, $y_{\pi_i}$~--- истинная релевантность новости на $i$-м месте, а $\mathrm{IDCG}@K$~--- максимально возможный $\mathrm{DCG}@K$ при идеальном упорядочивании.

Помимо метрик точности, для каждой парадигмы фиксируются показатели вычислительной стоимости в расчёте на один показ: суммарное число обработанных входных и выходных токенов $T_{\mathrm{total}} = T_{\mathrm{in}} + T_{\mathrm{out}}$ и среднее время работы (латентность) $\tau$ в секундах.

% ============================================================
\section{Описание методов}

В рамках работы реализуются три парадигмы \textit{zero-shot} ранжирования с использованием двух типов моделей: закрытой коммерческой LLM (через API) и открытой локально развёртываемой модели. Во всех режимах профиль пользователя $H_u$ формируется из заголовков и категорий последних $m$ прочитанных им новостей. Общая схема взаимодействия с LLM в трёх режимах представлена на рисунке~\ref{fig:schemes}, а сводное сравнение характеристик приведено в таблице~\ref{tab:methods}.

\begin{figure}[htbp]
  \centering
  \fbox{\parbox[c]{0.92\columnwidth}{%
    \footnotesize
    \textbf{1. Pointwise:} $[H_u,\ c_i] \rightarrow \text{LLM} \rightarrow \hat{s}_i \in [0, 1]$\\[3pt]
    \textbf{2. Pairwise:} $[H_u,\ c_i,\ c_j] \rightarrow \text{LLM} \rightarrow c_i \succ c_j$\\[3pt]
    \textbf{3. Listwise:} $[H_u,\ c_1, \dots, c_N] \rightarrow \text{LLM} \rightarrow (\pi_1, \dots, \pi_N)$%
  }}
  \caption{Схема формирования запросов к LLM в режимах pointwise, pairwise и listwise}
  \label{fig:schemes}
\end{figure}

\begin{table}[htbp]
  \caption{Сравнение парадигм zero-shot ранжирования}
  \label{tab:methods}
  \centering
  \begin{tabular}{@{}lccc@{}}
    \toprule
    \textbf{Метод} & \textbf{Запросов} & \textbf{Контекст} & \textbf{Учёт списка} \\
    \midrule
    Pointwise & $O(N)$ & Малый & Нет \\
    Pairwise  & $O(N \log N)$ & Малый & Локальный \\
    Listwise  & $O(1)$ & Большой & Глобальный \\
    \bottomrule
  \end{tabular}
\end{table}

\subsection{Поэлементный подход (Pointwise)}

В парадигме \textit{pointwise}~\cite{zhang2025recommendation} оценка каждого кандидата $c_i \in C_u$ производится независимо. Языковой модели подаётся промпт, содержащий историю пользователя $H_u$ и описание одного кандидата $c_i$. Перед моделью ставится задача оценить вероятность интереса пользователя к данной новости по заданной шкале либо вывести логиты целевых токенов подтверждения. Подход требует ровно $O(N)$ независимых запросов к LLM на один показ из $N$ кандидатов. Запросы легко распараллеливаются и требуют минимального размера контекстного окна, однако модель не видит остальных кандидатов и может назначать одинаковые оценки нескольким новостям (\textit{score ties}), что снижает точность взаимной сортировки.

\subsection{Попарный подход (Pairwise)}

В парадигме \textit{pairwise}~\cite{qin2024pairwise} модель решает задачу бинарного выбора более релевантной новости из пары кандидатов $(c_i, c_j)$ с учётом истории $H_u$. Для устранения позиционного смещения (\textit{position bias}), при котором LLM склонны отдавать предпочтение первому документу в промпте, проверка пары может выполняться в обоих порядках: $(c_i, c_j)$ и $(c_j, c_i)$. Полное попарное сравнение всех кандидатов требует $O(N^2)$ обращений к модели, а применение турнирных алгоритмов сортировки (например, Heapsort или метода скользящего окна) снижает сложность до $O(N \log N)$ или $O(K \cdot N)$. Задача выбора из двух альтернатив существенно проще для открытых моделей умеренного размера, но сопряжена с повышенным расходом входных токенов из-за многократной передачи истории $H_u$.

\subsection{Списочный подход (Listwise) и каскадная схема}

В парадигме \textit{listwise}~\cite{sun2023chatgpt} в один промпт помещаются сразу история пользователя $H_u$ и пронумерованный список из $N$ кандидатов. От модели требуется сгенерировать строку с перестановкой идентификаторов в порядке убывания релевантности (например, \texttt{[3] > [1] > [5]}). При большом числе кандидатов применяется итеративное переранжирование скользящим окном размера $w$ с шагом $s$, а также исследуется двухэтапная последовательная комбинация \textit{pointwise}-отбора и финального \textit{listwise}-переранжирования топ-кандидатов для снижения вычислительных затрат при формировании новостной ленты.

\subsection{Базовые методы (Baselines)}

Для объективной проверки результатов LLM-ранжирования используются три базовых метода:
\begin{itemize}
  \item \textbf{BM25}~--- классический алгоритм лексического ранжирования на основе частот термов между склеенными заголовками истории $H_u$ и текстом кандидата $c_i$;
  \item \textbf{Эмбеддинги (Dense Retrieval)}~--- ранжирование по косинусной близости между усреднённым вектором истории пользователя и вектором новости-кандидата, полученными с помощью предобученного трансформерного энкодера;
  \item \textbf{NRMS}~\cite{wu2019nrms}~--- специализированная нейросетевая модель для рекомендации новостей на базе механизма многоголового самовнимания (\textit{multi-head self-attention}), обученная на тренировочной выборке датасета MIND.
\end{itemize}

% ============================================================
\begin{thebibliography}{99}

\bibitem{sun2023chatgpt} Sun W., Yan L., Ma X., Wang S., Ren P., Chen Z., Yin D., Ren Z. Is ChatGPT Good at Search? Investigating Large Language Models as Re-Ranking Agents // \textit{Proceedings of the 2023 Conference on Empirical Methods in Natural Language Processing (EMNLP 2023)}. — 2023. — P.~14918–14937. — URL: \url{https://doi.org/10.18653/v1/2023.emnlp-main.923}

\bibitem{zhang2025recommendation} Zhang J., Xie R., Hou Y., Zhao W. X., Lin L., Wen J.-R. Recommendation as Instruction Following: A Large Language Model Empowered Recommendation Approach // \textit{ACM Transactions on Information Systems}. — 2025. — Vol.~43, №~5. — P.~1–37. — URL: \url{https://doi.org/10.1145/3708882}

\bibitem{qin2024pairwise} Qin Z., Jagerman R., Hui K., Zhuang H., Wu J., Yan L., Shen J., Liu T., Liu J., Metzler D., Wang X., Bendersky M. Large Language Models are Effective Text Rankers with Pairwise Ranking Prompting // \textit{Findings of the Association for Computational Linguistics: NAACL 2024}. — 2024. — P.~1504–1518. — URL: \url{https://doi.org/10.18653/v1/2024.findings-naacl.97}

\bibitem{wu2020mind} Wu F., Qiao Y., Chen J.-H., Wu C., Qi T., Lian J., Liu D., Xie X., Gao J., Wu W., Zhou M. MIND: A Large-scale Dataset for News Recommendation // \textit{Proceedings of the 58th Annual Meeting of the Association for Computational Linguistics (ACL 2020)}. — 2020. — P.~3597–3606. — URL: \url{https://doi.org/10.18653/v1/2020.acl-main.331}

\bibitem{wu2019nrms} Wu C., Wu F., Ge S., Qi T., Huang Y., Xie X. Neural News Recommendation with Multi-Head Self-Attention // \textit{Proceedings of the 2019 Conference on Empirical Methods in Natural Language Processing (EMNLP-IJCNLP 2019)}. — 2019. — P.~6389–6394. — URL: \url{https://doi.org/10.18653/v1/D19-1671}

\bibitem{sosedka2024telegram} Соседка А. В., Ибрагимов П. И. Количественный и качественный анализ аудитории Telegram в разрезе рекомендаций с использованием больших языковых моделей // \textit{Сборник статей МИСИС (весенний)}. — 2024. — С.~117–121.

\bibitem{beresnev2023chatbots} Береснев Д. В. Анализ подходов к использованию предобученных моделей в разработке корпоративных чат-ботов // \textit{Сборник статей МИСИС (осенний)}. — 2023. — С.~70–79.

\end{thebibliography}

\end{document}
